% =====================================================================
\chapter{Security and performance analysis in the paper}\label{ch:p1analysis}
% =====================================================================
\section{Security analysis in the paper}
Section~V of the paper gives five theorems. We summarise each one and how the paper argues it.
\begin{center}
\small
\begin{tabularx}{\linewidth}{@{}p{3.8cm} X@{}}
\toprule
\textbf{Theorem (paper)} & \textbf{What it claims and how it is argued}\\
\midrule
1 -- Message unforgeability & secure against chosen-message attacks in the random-oracle model if CDH is hard; the proof embeds the CDH instance
$(s_{tu}P,\ (\sum_i\upsilon_i\xi_i)P)$ and uses two forgeries with different $h_2$ values\\
2 -- User anonymity and unlinkability & temporary identities $\ID_i=h_3(id_i,\ts_2,\xi_i)$ change every session, so eavesdroppers cannot link messages\\
3 -- Replay attack resistance & every credential includes a fresh timestamp, so old messages fail the current check\\
4 -- Conditional privacy & real identities stay hidden, but the law-enforcement agency can reveal a malicious UAV when the TA reports it\\
5 -- Session key establishment & every verified UAV gets the same group key $\kappa_{tu}=\Lambda(\sigma_i)\bmod\sigma_i$; remaining UAVs can decrypt
updated keys ``if the $\sigma_i$ has not been removed by TUAV''\\
\bottomrule
\end{tabularx}
\end{center}
In its Table~II the paper compares itself with four earlier schemes on eleven properties (F1--F11) and ticks all eleven for itself, including
\emph{key escrow resilience}, \emph{forward secrecy} and \emph{non-repudiation}.

\section{Performance analysis in the paper}
The paper measures basic operations with the MIRACL library at an 80-bit security level: a supersingular curve with a 512-bit field for the pairing
group $\G_1$, and an ordinary 160-bit curve for the group $\G$.
\begin{center}
\small
\begin{tabular}{@{}l l r@{}}
\toprule
\textbf{Symbol} & \textbf{Operation} & \textbf{Time (ms)}\\
\midrule
$T_P$ & bilinear pairing & 7.2232\\
$T_{SM\cdot P}$ & scalar multiplication in $\G_1$ & 2.6784\\
$T_{H\cdot MP\cdot P}$ & MapToPoint hash & 8.1940\\
$T_{SM\cdot E}$ & scalar multiplication on the ordinary curve & 0.9018\\
$T_{PA\cdot E}$ & point addition on the ordinary curve & 0.0107\\
$T_{ME}$ & modular exponentiation & 0.6256\\
$T_{SH}$ & hash & 0.0012\\
\bottomrule
\end{tabular}
\end{center}
With these constants the paper reports: UAV signing (US) $2T_{SM\cdot E}+T_{PA\cdot E}+5T_{SH}\approx1.82$\,ms, single verification (SA)
$\approx2.16$\,ms, and group authentication (GA) $1.2524n+0.9054$\,ms, i.e.\ about $151$\,ms for $n=120$ UAVs. For GA the paper notes that the pairing
work is done by the TA ``prior to the verification of TUAV'', so it counts the TUAV side.

\begin{figure}[H]
\centering
\begin{tikzpicture}
\begin{axis}[width=0.88\linewidth,height=6.6cm,xmin=0,xmax=120,ymin=0,ymax=2100,xlabel={Number of UAVs $n$},ylabel={Group-authentication time (ms)},
  xtick={0,20,...,120},grid=major,grid style={black!10},legend cell align=left,legend style={font=\footnotesize,at={(0.02,0.98)},anchor=north west,draw=black!40},
  tick label style={font=\small},label style={font=\small},mark repeat=2,scaled y ticks=false,y tick label style={/pgf/number format/.cd,fixed,1000 sep={}}]
  \addplot[thick,color=plotB,mark=square*,mark size=1.6pt,domain=0:120,samples=13] {16.2292*x+37.0868}; \addlegendentry{Kumar et al.: $16.23n+37.09$}
  \addplot[thick,color=plotD,mark=triangle*,mark size=2pt,domain=0:120,samples=13] {13.6534*x+29.7952}; \addlegendentry{Xu et al.: $13.65n+29.80$}
  \addplot[thick,color=plotC,mark=diamond*,mark size=2pt,domain=0:120,samples=13] {5.4252*x+45.2124}; \addlegendentry{Mei et al.: $5.43n+45.21$}
  \addplot[thick,color=plotE,mark=o,mark size=1.6pt,domain=0:120,samples=13] {2.7126*x+7.2232}; \addlegendentry{Ali et al.: $2.71n+7.22$}
  \addplot[very thick,color=plotA,mark=*,mark size=1.6pt,domain=0:120,samples=13] {1.2524*x+0.9054}; \addlegendentry{Base paper (TUAV side): $1.25n+0.91$}
  \node[font=\footnotesize,anchor=south east] at (axis cs:119,165) {151 ms at $n=120$};
\end{axis}
\end{tikzpicture}
\caption{Group-authentication (GA) time against the number of UAVs, plotted from the cost formulas in Section~VI-A of the base paper for the four schemes
it compares with (Kumar et al., Xu et al., Mei et al., Ali et al.) and for itself. The base paper's line counts only the TUAV side; the TA's pairings
are not included.}
\label{fig:p1cost}
\end{figure}

For communication, the paper uses 128-byte elements in $\G_1$, 40-byte elements in $\G$, 20-byte hashes and 4-byte timestamps, and reports lower
overhead than the four compared schemes (its Table~V).

\section{Open questions we took forward}
While studying the paper we noted several points that we wanted to check carefully:
\begin{enumerate}[leftmargin=*]
\item The TA stores every UAV's secret $s_i$ (IS2), yet the paper claims key-escrow resilience and non-repudiation. \hfill $\rightarrow$ F1
\item The coefficients of $\Lambda(x)$ are sent in the clear and only the constant term is masked. \hfill $\rightarrow$ F2
\item The decrypting key $\sigma_i=h_1(s_i\xi_i)P$ uses no fresh secret, yet forward secrecy is claimed. \hfill $\rightarrow$ F3
\item Every login goes through the TA over the satellite, and only one TUAV is assumed for the whole region (IS4). \hfill $\rightarrow$ C1, C2
\item Nothing is said about a TUAV that is hacked or physically captured. \hfill $\rightarrow$ C3, C4
\end{enumerate}
